\documentclass{beamer}
\usetheme{Madrid}
\usecolortheme{beaver}
\usepackage{graphicx}
\usepackage{listings}
\usepackage{amsmath}
\usepackage{hyperref}
\usepackage{xcolor}

\definecolor{codegreen}{rgb}{0,0.6,0}
\definecolor{codegray}{rgb}{0.5,0.5,0.5}
\definecolor{codepurple}{rgb}{0.58,0,0.82}
\definecolor{backcolour}{rgb}{0.95,0.95,0.92}

\lstdefinestyle{mystyle}{
    backgroundcolor=\color{backcolour},   
    commentstyle=\color{codegreen},
    keywordstyle=\color{magenta},
    numberstyle=\tiny\color{codegray},
    stringstyle=\color{codepurple},
    basicstyle=\ttfamily\footnotesize,
    breakatwhitespace=false,         
    breaklines=true,                 
    captionpos=b,                    
    keepspaces=true,                 
    numbers=left,                    
    numbersep=5pt,                  
    showspaces=false,                
    showstringspaces=false,
    showtabs=false,                  
    tabsize=2
}

\lstset{style=mystyle}

\title{Traffic Signs Recognition using CNN in Keras}
\author{Presentation Based on Deepak's GitHub Project}
\date{\today}

\begin{document}

\begin{frame}
\titlepage
\end{frame}

\begin{frame}{Table of Contents}
\tableofcontents
\end{frame}

\section{Introduction}

\begin{frame}{Introduction}
\begin{itemize}
    Traffic sign recognition is a critical component for autonomous vehicles and driver assistance systems
    Enables safer navigation by interpreting road signs automatically
    Deep learning approaches, particularly Convolutional Neural Networks (CNNs), have shown remarkable accuracy in this task
    This presentation covers a comprehensive implementation of traffic sign recognition using CNN with Keras
\end{itemize}
\end{frame}

\begin{frame}{Project Overview}
\begin{itemize}
    {Dataset}: German Traffic Sign Recognition Benchmark (GTSRB)
    {Classification}: 43 different classes of traffic signs
    {Framework}: TensorFlow and Keras
    {Approach}: CNN architecture for feature extraction and classification
    {Implementation}: Complete pipeline from data preprocessing to web application deployment
\end{itemize}
\end{frame}

\section{Dataset Analysis}

\begin{frame}{Dataset Description}
\begin{itemize}
    The German Traffic Sign Recognition Benchmark (GTSRB)
    Contains more than 50,000 images of traffic signs
    43 classes of signs (speed limits, no entry, stop, etc.)
    Images are of varying sizes and lighting conditions
    Real-world challenges: occlusion, weather effects, and viewpoint variations
    Data organized into training, validation, and test sets
\end{itemize}
\end{frame}

\begin{frame}{Exploratory Data Analysis}
\begin{lstlisting}[language=Python]
# From eda.py
# Class distribution plot
plt.figure(figsize=(12, 6))
sns.barplot(x=class_labels, y=class_counts)
plt.title("Class Distribution of Traffic Signs")
plt.xlabel("Class")
plt.ylabel("Number of Samples")
plt.xticks(class_labels)
plt.savefig(os.path.join(output_dir, 'class_distribution.png'))
\end{lstlisting}

\begin{itemize}
    Analysis of class distribution shows imbalanced representation
    Visualization of sample images from each class
    Data insights inform preprocessing decisions
\end{itemize}
\end{frame}

\begin{frame}{Sample Images Visualization}
\begin{lstlisting}[language=Python]
# From eda.py
# Plot sample images
plt.figure(figsize=(15, 10))
for i in range(1, 11):
    plt.subplot(2, 5, i)
    plt.imshow(sample_images[i - 1])
    plt.title(f"Class: {sample_labels[i - 1]}")
    plt.axis('off')
plt.tight_layout()
\end{lstlisting}

\begin{itemize}
Visualization helps understand diversity in the dataset
Variety in illumination, perspective, and quality
Challenges for model training become apparent
\end{itemize}
\end{frame}

\section{Data Preprocessing}

\begin{frame}{Data Preprocessing Strategy}
\begin{itemize}
    {Image Resizing}: All images resized to 30×30 pixels for uniform input
    {Normalization}: Pixel values scaled to [0,1] range
    {Label Encoding}: One-hot encoding for multi-class classification
    {Train-Test Split}: Data divided for training and evaluation
    {Data Augmentation}: (Optional) To address class imbalance
\end{itemize}
\end{frame}

\begin{frame}{Preprocessing Implementation}
\begin{lstlisting}[language=Python]
# Data preprocessing pipeline
# Convert images to numpy arrays and normalize
X_train = np.array(images) / 255.0

# One-hot encode the labels
y_train = np.zeros((n_samples, n_classes))
for i, label in enumerate(labels):
    y_train[i, label] = 1.0

# Save the preprocessed data
np.save(os.path.join(output_dir, 'X_train.npy'), X_train)
np.save(os.path.join(output_dir, 'y_train.npy'), y_train)
\end{lstlisting}

\begin{itemize}
    Preprocessed data saved as NumPy arrays for efficient loading
    Consistent preprocessing ensures reliable model performance
\end{itemize}
\end{frame}

\section{CNN Architecture}

\begin{frame}{CNN Architecture Overview}
\begin{itemize}
    Convolutional Neural Networks are ideal for image classification
    Architecture components:
    \begin{itemize}
        Convolutional layers: Feature extraction
        Pooling layers: Dimensionality reduction
        Dropout layers: Regularization to prevent overfitting
        Dense layers: Classification based on extracted features
    \end{itemize}
    Multi-stage architecture for hierarchical feature learning
\end{itemize}
\end{frame}

\begin{frame}{Model Architecture Implementation}
\begin{lstlisting}[language=Python]
# From model_training.py
model = Sequential()
# First Layer
model.add(Conv2D(filters=32, kernel_size=(5, 5), 
                 activation='relu', input_shape=X_train.shape[1:]))
model.add(Conv2D(filters=32, kernel_size=(5, 5), activation='relu'))
model.add(MaxPooling2D(pool_size=(2, 2)))
model.add(Dropout(rate=0.25))
# Second Layer 
model.add(Conv2D(filters=64, kernel_size=(3, 3), activation='relu'))
model.add(Conv2D(filters=64, kernel_size=(3, 3), activation='relu'))
model.add(MaxPooling2D(pool_size=(2, 2)))
model.add(Dropout(rate=0.25))
# Dense Layer
model.add(Flatten())
model.add(Dense(256, activation='relu'))
model.add(Dropout(rate=0.5))
model.add(Dense(43, activation='softmax'))
\end{lstlisting}
\end{frame}

\begin{frame}{Architectural Design Choices}
\begin{itemize}
    {Convolutional Layers}:
    \begin{itemize}
        Two blocks with increasing filter sizes (32 → 64)
        Smaller kernel sizes in deeper layers (5×5 → 3×3)
        ReLU activation for non-linearity
    \end{itemize}
    {Regularization Strategy}:
    \begin{itemize}
        Dropout layers with varying rates (0.25 in conv blocks, 0.5 in dense)
        MaxPooling for spatial downsampling and parameter reduction
    \end{itemize}
    {Output Layer}:
    \begin{itemize}
        43 neurons with softmax activation for multi-class probabilities
    \end{itemize}
\end{itemize}
\end{frame}

\section{Model Training}

\begin{frame}{Training Configuration}
\begin{lstlisting}[language=Python]
# From model_training.py
model.compile(loss='categorical_crossentropy', 
              optimizer='adam', 
              metrics=['accuracy'])

# Train the model
history = model.fit(X_train, y_train, 
                    batch_size=args.batch_size, 
                    epochs=args.epochs, 
                    validation_data=(X_test, y_test))
\end{lstlisting}

\begin{itemize}
{Loss Function}: Categorical Cross-Entropy (appropriate for multi-class)
{Optimizer}: Adam (adaptive learning rate)
{Metrics}: Accuracy to monitor performance
{Batch Size}: 64 (default)
{Epochs}: 20 (default, configurable via command line)
\end{itemize}
\end{frame}

\begin{frame}{Model Training Workflow}
\begin{itemize}
    Complete training pipeline:
    \begin{enumerate}
        Load preprocessed data from NumPy arrays
        Build the CNN model
        Configure training parameters
        Train with validation monitoring
        Save the best model and training history
    \end{enumerate}
    Command-line interface for flexible configuration:
    \begin{itemize}
        Data directory
        Model save location
        Number of epochs
        Batch size
    \end{itemize}
\end{itemize}
\end{frame}

\section{Model Evaluation}

\begin{frame}{Evaluation Metrics}
\begin{itemize}
{Loss}: How far predictions deviate from actual labels
{Accuracy}: Percentage of correctly classified signs
{Training vs. Validation}: Monitoring for overfitting
{Confusion Matrix}: (Not implemented but valuable addition)
{Class-specific Performance}: Understanding per-class accuracy
\end{itemize}
\end{frame}

\begin{frame}{Evaluation Implementation}
\begin{lstlisting}[language=Python]
# From evaluation.py
# Load preprocessed data
X_test = np.load(os.path.join(data_dir, 'X_test.npy'))
y_test = np.load(os.path.join(data_dir, 'y_test.npy'))

# Load model
model = tf.keras.models.load_model(model_path)

# Evaluate model
loss, accuracy = model.evaluate(X_test, y_test)
print(f"Test Loss: {loss:.4f}, Test Accuracy: {accuracy * 100:.2f}%")

# Plot accuracy and loss
plt.figure(figsize=(12, 5))
plt.subplot(1, 2, 1)
plt.plot(history['accuracy'], label='Train Accuracy')
plt.plot(history['val_accuracy'], label='Validation Accuracy')
\end{lstlisting}
\end{frame}

\begin{frame}{Performance Visualization}
\begin{itemize}
    Training and validation curves provide insights:
    \begin{itemize}
        Convergence: Is the model learning?
        Overfitting: Is there a gap between training and validation?
        Early stopping points: When did improvement plateau?
    \end{itemize}
    Visual assessment complements numerical metrics
    Helps in refining the model architecture or training process
\end{itemize}
\end{frame}

\section{Inference}

\begin{frame}{Inference Pipeline}
\begin{lstlisting}[language=Python]
# From inference.py
# Load and preprocess the image
image = Image.open(image_path).resize((30, 30))
image_array = np.array(image) / 255.0
image_array = np.expand_dims(image_array, axis=0)

# Predict the class
prediction = model.predict(image_array)
predicted_class = np.argmax(prediction, axis=1)[0]
\end{lstlisting}

\begin{itemize}
    Standalone inference module for testing single images
    Same preprocessing steps ensure consistency
    Command-line interface for easy testing:
    \begin{itemize}
        Model path selection
        Test image specification
    \end{itemize}
\end{itemize}
\end{frame}

\section{Web Application}

\begin{frame}{Streamlit Web Application}
\begin{lstlisting}[language=Python]
# From streamlit_app.py
def run(args):
    model_path = args.model_path
    st.title("Traffic Sign Recognition App")
    # Load the model
    model = tf.keras.models.load_model(model_path)
    # File uploader
    uploaded_file = st.file_uploader("Choose a traffic sign image...", 
                                     type=["jpg", "png", "jpeg"])
    if uploaded_file is not None:
        # Display the uploaded image
        image = Image.open(uploaded_file).resize((30, 30))
        st.image(image, caption='Uploaded Image', use_column_width=True)
        
        # Preprocess and predict
        image_array = np.array(image) / 255.0
        image_array = np.expand_dims(image_array, axis=0)
        prediction = model.predict(image_array)
        predicted_class = np.argmax(prediction, axis=1)[0]
        st.write(f"Predicted Class: {predicted_class}")
\end{lstlisting}
\end{frame}

\begin{frame}{User Interface Features}
\begin{itemize}
    Interactive web application built with Streamlit
    Features:
    \begin{itemize}
        Image upload capability
        Real-time preprocessing and inference
        Immediate display of results
        User-friendly interface
    \end{itemize}
    Easily deployable for demonstrations
    Potential enhancements:
    \begin{itemize}
        Class descriptions instead of just numbers
        Confidence scores
        Multiple prediction options
    \end{itemize}
\end{itemize}
\end{frame}

\section{Project Pipeline}

\begin{frame}{Complete Project Pipeline}
\begin{itemize}
    {Data Collection}: Using GTSRB dataset
    {Exploratory Data Analysis}: Understanding data characteristics
    {Data Preprocessing}: Resizing, normalization, and encoding
    {Model Development}: CNN architecture design
    {Model Training}: Parameter optimization and validation
    {Model Evaluation}: Performance assessment
    {Inference}: Real-world testing
    {Application Development}: User interface creation
\end{itemize}
\end{frame}

\begin{frame}{Command Line Execution}
\begin{lstlisting}[language=bash]
# Run exploratory data analysis
python eda.py --data_dir data --output_dir outputs

# Run model training
python model_training.py --data_dir data --model_dir models 
                        --epochs 25 --batch_size 64

# Evaluate the model
python evaluation.py --data_dir data --model_dir models 
                    --output_dir outputs

# Run inference on a single image
python inference.py --model_dir models 
                   --image_path path/to/test_image.jpg

# Launch the Streamlit web app
streamlit run streamlit_app.py -- --model_path models/best_model.h5
\end{lstlisting}
\end{frame}

\section{Future Improvements}

\begin{frame}{Potential Enhancements}
\begin{itemize}
    {Model Architecture}:
    \begin{itemize}
        Transfer learning with pre-trained networks (ResNet, MobileNet)
        Attention mechanisms for focusing on critical sign features
    \end{itemize}
    {Data Handling}:
    \begin{itemize}
        Advanced augmentation techniques for better generalization
        Class balancing to improve performance on underrepresented signs
    \end{itemize}
    {Application}:
    \begin{itemize}
        Real-time video processing for continuous sign detection
        Integration with geolocation for context-aware predictions
        Multi-language support for sign descriptions
    \end{itemize}
\end{itemize}
\end{frame}

\begin{frame}{Conclusion}
\begin{itemize}
    Traffic sign recognition is crucial for autonomous driving safety
    CNN-based approach provides high accuracy and robust performance
    Complete pipeline from data preparation to deployment demonstrated
    The modular design allows for easy extension and improvement
    Future work can focus on model optimization and real-world integration
\end{itemize}
\end{frame}

\begin{frame}{References}
\begin{itemize}
    GitHub Repository: 
    \url{https://github.com/deepak2233/Traffic-Signs-Recognition-using-CNN-Keras}
    German Traffic Sign Recognition Benchmark:
    \url{https://benchmark.ini.rub.de/gtsrb_news.html}
    Krizhevsky, A., Sutskever, I., \& Hinton, G. E. (2012). ImageNet classification with deep convolutional neural networks. Advances in neural information processing systems, 25.
    Sermanet, P., \& LeCun, Y. (2011, July). Traffic sign recognition with multi-scale convolutional networks. In The 2011 International Joint Conference on Neural Networks (pp. 2809-2813). IEEE.
\end{itemize}
\end{frame}

\end{document}